\section{Phase 3b: Architectural Interventions and Agent-Type \texorpdfstring{$\Delta\theta$}{Delta-theta}}
\label{sec:phase3b}

\subsection{Motivation}

Phase 3a established that system-prompt strictness primarily unlocks
citation \emph{format} rather than compliance \emph{reasoning}---with
only the largest model showing a genuine reasoning gain (+2.5 logits
under the lenient judge). Phase 3b tests whether \emph{structural}
changes to the agent---knowledge retrieval, self-critique, and step
decomposition---produce measurable $\Delta\theta$ beyond what prompt
engineering delivers, using the same frozen 1,284-item bank as the
measuring instrument.

\subsection{Agent Variants}

We implement three agent architectures, each a drop-in replacement for
the zero-shot baseline in the Arena runner:

\paragraph{RetrievalAgent.}
Before the single LLM call, the agent embeds the item's context and
question with the \texttt{all-MiniLM-L6-v2} sentence-transformer model
(22\,MB, local, no API call) and retrieves the top-$k$ ($k=3$) CFR
sections by cosine similarity from a pre-built numpy index covering
CFR Parts 11, 50, 58, 211 (Title 21) and Part 46 (Title 45).  The
retrieved sections are prepended to the system prompt as a
\emph{Relevant Regulations} block; all other parameters are identical
to the zero-shot baseline (same strictness level, model, temperature).
This isolates the effect of in-context regulatory grounding without
changing the agent's reasoning process.

\paragraph{CriticAgent.}
The agent makes an initial zero-shot answer, then issues a critique
prompt---``Review your previous answer for citation accuracy and
compliance determination. Correct any errors.''---and finally a
synthesis step.  With one critique round this costs three Groq calls
per item (default).  All turns share the same conversation thread so
the model has access to its own reasoning history.

\paragraph{StepDecompositionAgent.}
The agent breaks the compliance analysis into four sequential Groq
calls, each building on the prior step's output: (1)~identify the
regulated activity and parties; (2)~identify the governing CFR/ICH
provision; (3)~assess compliance or violation; (4)~state the final
determination with section-level citation.  This mirrors a structured
legal-analysis workflow and tests whether explicit chain-of-thought
scaffolding improves compliance reasoning.

\subsection{Experimental Grid}

The Phase 3b grid crosses two models (llama-3.3-70b-versatile,
llama-3.1-8b-instant) $\times$ one temperature (0.4) $\times$ two
strictness levels (none, strict) $\times$ four agent types (zero\_shot,
retrieval, critic, step\_decomp), yielding 16 examinee cells.  Every
cell is administered the full frozen 1,284-item bank, anchored to the
Phase 1 baseline ($\theta = 0$) using the same re-anchoring procedure
as Phase 3a (Section~\ref{sec:delta_theta}).

\subsection{Results}

\begin{table}[h]
\centering
\caption{Phase 3b $\Delta\theta$ by agent type relative to zero-shot baseline
         (same model, temperature, strictness). \todo{Fill after phase3b run.}}
\label{tab:phase3b}
\begin{tabular}{llcc}
\toprule
Agent type & Model & $\Delta\theta$ (strictness=none) & $\Delta\theta$ (strictness=strict) \\
\midrule
retrieval   & llama-3.3-70b & \todo{} & \todo{} \\
critic      & llama-3.3-70b & \todo{} & \todo{} \\
step\_decomp & llama-3.3-70b & \todo{} & \todo{} \\
retrieval   & llama-3.1-8b  & \todo{} & \todo{} \\
critic      & llama-3.1-8b  & \todo{} & \todo{} \\
step\_decomp & llama-3.1-8b  & \todo{} & \todo{} \\
\bottomrule
\end{tabular}
\end{table}

\todo{Interpret Table~\ref{tab:phase3b}: does in-context CFR retrieval
move $\theta$ beyond the zero-shot baseline? Is the retrieval gain
larger under strict (citation-requiring) than lenient grading---which
would suggest the agent is using the retrieved sections for format
rather than reasoning? Does step decomposition produce a consistent
reasoning gain across model sizes?}
